La formulación de los requisitos de calidad plantea una dificultad mayor que la descripción de las funciones de búsqueda y gestión de listas: exige transformar expresiones generales en obligaciones medibles. Un ejemplo es pasar de «la búsqueda será muy rápida» a una condición con número de consultas, límite de tiempo y entorno de medición (R-11). Aunque la segunda versión es más verificable, el umbral de tres segundos todavía debe justificarse mediante necesidades de usuarios o criterios del proyecto. Lo mismo ocurre con la privacidad: la presencia de «Sesión privada» en el menú permite identificar la función, pero no comprobar que la actividad quede oculta (Figura~\ref{fig:privada}). Para demostrarlo haría falta un protocolo con cuentas autorizadas y registros de resultados.

El ejercicio diferencia tres niveles de evidencia que no conviene confundir: una captura muestra un estado de la interfaz, una prueba funcional comprueba una operación y una prueba de calidad mide una propiedad bajo condiciones definidas. Además, los requisitos se deben revisar como conjunto para evitar conflictos, por ejemplo entre permisos de las cuentas Free y Premium. El principal aprendizaje técnico consiste en no confundir un requisito redactado con precisión con uno ya validado o cumplido. Esta reflexión describe hallazgos del análisis documental; el estudiante debe revisarla y agregar sus propias dificultades antes de atribuírsela como experiencia personal.
